% Unit 7 exam -- 30 MCQ + 1 short + 1 long FRQ = 44 pts, 76 min.
\documentclass{apchem-exam}

\apcunit{7}
\apcunittitle{Equilibrium}
\apcdoctitle{Unit Exam}
\apcsubtitle{CED 7.1--7.12 \textbullet\ Mon Jan 25 \textbullet\ 76 minutes \textbullet\ 44 points}

\begin{document}
\apctitleblock

\begin{apcnote}[Format]
Section I: 30 multiple choice, \textbf{no calculator}, 45 minutes.
Section II: one short and one long free response, calculator permitted,
31 minutes. Constants are provided with Section II.
\end{apcnote}

\examsection{I}{Multiple Choice}{30 questions \textbullet\ 45 minutes \textbullet\ 1 point each}{\nocalculator}

% ---- dynamic equilibrium ----------------------------------------------------
\begin{mcq}
  At equilibrium,
  \choices
    \correct{the forward and reverse rates are equal}
    \choice{the concentrations of reactants and products are equal}
    \choice{all reaction has ceased}
    \choice{$Q$ is zero}
\end{mcq}
\why{Equal rates, not equal amounts, and nothing has stopped.}
\distractor{B}{the single most common misconception in the unit}

\begin{mcq}
  On a concentration--time graph, equilibrium is reached when
  \choices
    \correct{both curves become horizontal}
    \choice{the two curves cross}
    \choice{the product curve reaches its maximum possible value}
    \choice{the reactant curve reaches zero}
\end{mcq}
\why{Levelling off means concentrations have stopped changing.}

\begin{mcq}
  Starting from pure reactants, the reverse rate initially is
  \choices
    \correct{zero, and increases as product forms}
    \choice{at its maximum}
    \choice{equal to the forward rate}
    \choice{constant throughout}
\end{mcq}
\why{No product means no reverse reaction yet.}

% ---- writing K --------------------------------------------------------------
\begin{mcq}
  For $\ce{2SO2(g) + O2(g) <=> 2SO3(g)}$, $K$ equals
  \choices
    \correct{$[\ce{SO3}]^2/([\ce{SO2}]^2[\ce{O2}])$}
    \choice{$[\ce{SO3}]/([\ce{SO2}][\ce{O2}])$}
    \choice{$([\ce{SO2}]^2[\ce{O2}])/[\ce{SO3}]^2$}
    \choice{$[\ce{SO3}]^2/([\ce{SO2}][\ce{O2}])$}
\end{mcq}
\why{Products over reactants, each raised to its coefficient.}

\begin{mcq}
  For $\ce{CaCO3(s) <=> CaO(s) + CO2(g)}$, $K$ equals
  \choices
    \choice{$[\ce{CaO}][\ce{CO2}]/[\ce{CaCO3}]$}
    \correct{$[\ce{CO2}]$}
    \choice{$[\ce{CaO}][\ce{CO2}]$}
    \choice{$1/[\ce{CO2}]$}
\end{mcq}
\why{Both solids are omitted.}

\begin{mcq}
  Water appears in the equilibrium expression when it is
  \choices
    \correct{a gas or a dissolved species, but not the pure liquid solvent}
    \choice{always}
    \choice{never}
    \choice{only when it is a product}
\end{mcq}
\why{The state symbol decides.}

\begin{mcq}
  Pure solids are omitted from $K$ because
  \choices
    \correct{their concentration does not change as the reaction proceeds}
    \choice{they do not react}
    \choice{their concentration is zero}
    \choice{they are always in excess}
\end{mcq}
\why{A larger lump is not more concentrated.}

% ---- Q vs K -----------------------------------------------------------------
\begin{mcq}
  If $Q < K$, the reaction
  \choices
    \correct{proceeds to the right}
    \choice{proceeds to the left}
    \choice{is at equilibrium}
    \choice{stops}
\end{mcq}
\why{Too few products relative to equilibrium.}

\begin{mcq}
  $Q$ differs from $K$ in that
  \choices
    \correct{$Q$ uses concentrations at any moment, $K$ uses equilibrium
      ones}
    \choice{$Q$ omits solids while $K$ includes them}
    \choice{$Q$ applies only to gases}
    \choice{$Q$ is the reciprocal of $K$}
\end{mcq}
\why{The algebraic form is identical.}

\begin{mcq}
  For $\ce{H2 + I2 <=> 2HI}$ with $K = 4.5$, a mixture with
  $[\ce{H2}] = [\ce{I2}] = [\ce{HI}] = 0.50$~M will
  \choices
    \correct{shift right}
    \choice{shift left}
    \choice{be at equilibrium}
    \choice{cannot be determined}
\end{mcq}
\why{$Q = 1.0 < 4.5$.}

\begin{mcq}
  Which tool works for a mixture that was never at equilibrium?
  \choices
    \correct{comparing $Q$ with $K$}
    \choice{Le Ch\^atelier's principle}
    \choice{neither}
    \choice{both equally}
\end{mcq}
\why{Le Ch\^atelier describes disturbing an existing equilibrium.}

% ---- magnitude and properties ----------------------------------------------
\begin{mcq}
  A reaction with $K = 3\times10^{-8}$
  \choices
    \correct{strongly favors reactants}
    \choice{strongly favors products}
    \choice{has appreciable amounts of both}
    \choice{proceeds very slowly}
\end{mcq}
\why{A tiny $K$ says the equilibrium lies far left --- it says nothing
about rate.}
\distractor{D}{$K$ and rate are independent}

\begin{mcq}
  If $K = 4.0$ for a reaction, $K$ for its reverse is
  \choices
    \choice{$-4.0$}
    \correct{0.25}
    \choice{16}
    \choice{2.0}
\end{mcq}
\why{$1/K$.}

\begin{mcq}
  If all coefficients in an equation are doubled, $K$ becomes
  \choices
    \choice{$2K$}
    \correct{$K^2$}
    \choice{$K/2$}
    \choice{unchanged}
\end{mcq}
\why{Each concentration term is squared.}

\begin{mcq}
  Two reactions with constants $K_1$ and $K_2$ are added. The overall $K$
  is
  \choices
    \correct{$K_1K_2$}
    \choice{$K_1 + K_2$}
    \choice{$K_1/K_2$}
    \choice{$K_1 - K_2$}
\end{mcq}
\why{Constants multiply where enthalpies add.}

\begin{mcq}
  $K_p$ and $K_c$ are numerically equal when
  \choices
    \correct{$\Delta n = 0$}
    \choice{the temperature is \SI{273}{\kelvin}}
    \choice{$K = 1$}
    \choice{all species are gases}
\end{mcq}
\why{$K_p = K_c(RT)^{\Delta n}$.}

% ---- ICE --------------------------------------------------------------------
\begin{mcq}
  In an ICE table for $\ce{A <=> 2B}$, if A changes by $-x$, then B changes
  by
  \choices
    \correct{$+2x$}
    \choice{$+x$}
    \choice{$-2x$}
    \choice{$+x/2$}
\end{mcq}
\why{The change row follows the coefficients.}

\begin{mcq}
  The small-$x$ approximation is justified when
  \choices
    \correct{$x$ is less than about 5\% of the initial concentration}
    \choice{$K$ is greater than 1}
    \choice{the reaction is exothermic}
    \choice{there is only one product}
\end{mcq}
\why{And the percentage must be reported, not assumed.}

\begin{mcq}
  Which situation makes the small-$x$ approximation most likely to fail?
  \choices
    \choice{$K = 10^{-10}$, $[\ce{A}]_0 = 1.0$~M}
    \choice{$K = 10^{-8}$, $[\ce{A}]_0 = 0.50$~M}
    \correct{$K = 1.0$, $[\ce{A}]_0 = 0.10$~M}
    \choice{$K = 10^{-6}$, $[\ce{A}]_0 = 2.0$~M}
\end{mcq}
\why{It is the ratio $[\ce{A}]_0/K$ that matters, and here it is only
0.1.}

\begin{mcq}
  For $\ce{H2 + I2 <=> 2HI}$ with equal starting concentrations, the
  equilibrium expression can be solved without the quadratic formula
  because
  \choices
    \correct{both sides are perfect squares}
    \choice{$K$ is large}
    \choice{$\Delta n = 0$}
    \choice{the reaction is exothermic}
\end{mcq}
\why{Taking the square root gives a linear equation.}

\begin{mcq}
  A negative root of an equilibrium quadratic is discarded because
  \choices
    \correct{it would imply a negative concentration}
    \choice{$K$ must be positive}
    \choice{the quadratic formula only gives one root}
    \choice{it corresponds to the reverse reaction}
\end{mcq}
\why{Only physically possible roots are kept.}

% ---- Le Chatelier -----------------------------------------------------------
\begin{mcq}
  Adding more reactant to a system at equilibrium
  \choices
    \correct{shifts it right and leaves $K$ unchanged}
    \choice{shifts it right and increases $K$}
    \choice{shifts it left and decreases $K$}
    \choice{has no effect}
\end{mcq}
\why{Only temperature changes $K$.}
\distractor{B}{the most-repeated error on Le Ch\^atelier questions}

\begin{mcq}
  Adding an inert gas at constant volume
  \choices
    \correct{causes no shift, because no partial pressure changes}
    \choice{shifts toward fewer moles of gas}
    \choice{shifts toward more moles of gas}
    \choice{increases $K$}
\end{mcq}
\why{Total pressure rises, but every partial pressure is unchanged.}

\begin{mcq}
  Compressing a gaseous equilibrium into a smaller volume shifts it
  \choices
    \correct{toward the side with fewer moles of gas}
    \choice{toward the side with more moles of gas}
    \choice{toward products always}
    \choice{not at all}
\end{mcq}
\why{The system opposes the increase in pressure.}

\begin{mcq}
  A catalyst added to a system at equilibrium
  \choices
    \correct{changes neither the position nor $K$}
    \choice{shifts the equilibrium right}
    \choice{increases $K$}
    \choice{shifts toward fewer moles of gas}
\end{mcq}
\why{It speeds both directions equally.}

\begin{mcq}
  For an exothermic forward reaction, raising the temperature
  \choices
    \correct{shifts it left and decreases $K$}
    \choice{shifts it right and increases $K$}
    \choice{shifts it left and leaves $K$ unchanged}
    \choice{has no effect}
\end{mcq}
\why{Heat behaves as a product; adding it pushes back toward reactants.}

\begin{mcq}
  The only disturbance that changes the value of $K$ is a change in
  \choices
    \correct{temperature}
    \choice{concentration}
    \choice{pressure}
    \choice{the amount of catalyst}
\end{mcq}
\why{Everything else changes $Q$, and the system returns to the same $K$.}

% ---- solubility -------------------------------------------------------------
\begin{mcq}
  For a salt \ce{MX2} with molar solubility $s$, $K_{sp}$ equals
  \choices
    \choice{$s^2$}
    \choice{$2s^3$}
    \correct{$4s^3$}
    \choice{$s^3$}
\end{mcq}
\why{$(s)(2s)^2$.}

\begin{mcq}
  Adding \ce{NaCl} to a saturated solution of \ce{AgCl} causes
  \choices
    \correct{the solubility to fall, with $K_{sp}$ unchanged}
    \choice{both the solubility and $K_{sp}$ to fall}
    \choice{the solubility to rise}
    \choice{no change}
\end{mcq}
\why{The common ion raises $Q$ above $K_{sp}$; the constant does not
move.}
\distractor{B}{$K_{sp}$ depends only on temperature}

\begin{mcq}
  $K_{sp}$ values may be compared directly to rank solubilities only when
  the salts
  \choices
    \correct{produce the same total number of ions}
    \choice{share a common cation}
    \choice{have similar molar masses}
    \choice{are all insoluble}
\end{mcq}
\why{Otherwise $s$ enters raised to different powers.}

\examsection{II}{Free Response}{2 questions \textbullet\ 31 minutes}{\calculatorok}

\begin{apcnote}[Data]
$K_p = K_c(RT)^{\Delta n}$ \textbullet\
$R = \SI{0.08206}{\litre\atm\per\mole\per\kelvin}$ \textbullet\
$K_{sp}$: \ce{AgCl} $1.6\times10^{-10}$ \textbullet\ \ce{CaF2}
$4.0\times10^{-11}$.
\end{apcnote}

% ---- FRQ 1 (short) ---------------------------------------------------------
\frq{short}{4}{CED 7.11--7.12 \textbullet\ solubility and the common ion}

\begin{parts}
  \item Write the dissolution equation and $K_{sp}$ expression for
        \ce{CaF2}. \answerlines[2]{\ce{CaF2(s) <=> Ca^2+(aq) + 2F^-(aq)};
        $K_{sp} = [\ce{Ca^2+}][\ce{F^-}]^2$}
  \item Calculate its molar solubility in pure water. \workspace[2.2cm]
        \keyonly{\begin{apcanswer}$4s^3 = 4.0\times10^{-11}$;
        $s^3 = 1.0\times10^{-11}$;
        $s = \mathbf{2.2\times10^{-4}}$~M\end{apcanswer}}
  \item Calculate its solubility in \SI{0.025}{\Molar} \ce{NaF}.
        \workspace[2.2cm]
        \keyonly{\begin{apcanswer}$[\ce{F-}] \approx 0.025$;
        $4.0\times10^{-11} = s(0.025)^2$;
        $s = \mathbf{6.4\times10^{-8}}$~M\end{apcanswer}}
  \item State what happened to $K_{sp}$ between (b) and (c) and why.
        \answerlines[2]{Nothing --- $K_{sp}$ depends only on temperature.
        Only the equilibrium position moved, toward less dissolved solid.}
\end{parts}

\begin{rubric}
  \rpoint{1}{(a) equation and expression, with fluoride squared}
  \rpoint{1}{(b) $s = 2.2\times10^{-4}$~M using $4s^3$}
  \rpoint{1}{(c) $s = 6.4\times10^{-8}$~M}
  \rpoint{1}{(d) $K_{sp}$ unchanged, attributed to temperature dependence
    --- ``$K_{sp}$ decreased'' earns 0}
\end{rubric}

% ---- FRQ 2 (long) ----------------------------------------------------------
\frq{long}{10}{CED 7.3, 7.7, 7.9, 7.10 \textbullet\ a complete equilibrium analysis}

At \SI{700}{\kelvin}, $K = 50.0$ for
$\ce{H2(g) + I2(g) <=> 2HI(g)}$. The forward reaction is exothermic.

\begin{parts}
  \item Write the equilibrium expression.
        \answerblank[6cm]{$[\ce{HI}]^2/([\ce{H2}][\ce{I2}])$}
  \item A flask is charged with \SI{1.00}{\Molar} \ce{H2} and
        \SI{1.00}{\Molar} \ce{I2}. Construct an ICE table and determine all
        three equilibrium concentrations. \workspace[3.6cm]
        \keyonly{\begin{apcanswer}Change: $-x$, $-x$, $+2x$.
        $K = (2x)^2/(1.00-x)^2 = 50.0$. Both sides are perfect squares, so
        $2x/(1.00-x) = 7.07$, giving $9.07x = 7.07$ and $x = 0.780$.
        $[\ce{H2}] = [\ce{I2}] = \mathbf{0.220}$~M;
        $[\ce{HI}] = \mathbf{1.56}$~M\end{apcanswer}}
  \item Verify your answer by substituting back into the expression.
        \workspace[1.8cm]
        \keyonly{\begin{apcanswer}$(1.56)^2/(0.220)^2 = 2.43/0.0484
        = \mathbf{50.2}$, which agrees with 50.0 to
        rounding\end{apcanswer}}
  \item A second flask at the same temperature contains
        $[\ce{H2}] = 0.40$, $[\ce{I2}] = 0.30$, $[\ce{HI}] = 0.40$~M.
        Calculate $Q$ and predict the direction of change, justifying your
        answer. \workspace[2.4cm]
        \keyonly{\begin{apcanswer}$Q = (0.40)^2/[(0.40)(0.30)] =
        \mathbf{1.3}$. Since $Q < K$, the mixture holds too few products
        and the reaction proceeds to the \textbf{right} until $Q$ reaches
        50.0.\end{apcanswer}}
  \item Explain why Le Ch\^atelier's principle is not the appropriate tool
        for part (d). \answerlines[2]{The second mixture was never at
        equilibrium, and Le Ch\^atelier describes only the response to
        disturbing an existing equilibrium. Comparing $Q$ with $K$ works
        for any mixture.}
  \item State the effect on the equilibrium position and on the value of
        $K$ of (i) adding more \ce{H2} and (ii) raising the temperature.
        \answerlines[3]{(i) The position shifts right; $K$ is unchanged,
        because adding \ce{H2} lowers $Q$ and the system simply returns to
        the same $K$. (ii) The position shifts left and $K$
        \emph{decreases}, because the forward reaction is exothermic, so
        heat behaves as a product.}
\end{parts}

\begin{rubric}
  \rpoint{1}{(a) correct expression}
  \rpoint{3}{(b) 1 pt ICE table with $+2x$ for \ce{HI}; 1 pt algebra to
    $x = 0.780$; 1 pt all three concentrations}
  \rpoint{1}{(c) substitution recovering $\approx 50$}
  \rpoint{2}{(d) 1 pt $Q = 1.3$; 1 pt direction \emph{with} the $Q < K$
    justification --- a bare ``shifts right'' earns 1}
  \rpoint{1}{(e) notes the mixture was never at equilibrium}
  \rpoint{2}{(f) 1 pt $K$ unchanged for the concentration change; 1 pt $K$
    decreases on heating, with the exothermic reason --- saying $K$
    increases earns 0}
\end{rubric}

\examtotal

\end{document}
